\documentclass[10pt]{article}
\usepackage[a4paper,margin=0.55in]{geometry}
\usepackage{xcolor}
\usepackage{multicol}
\usepackage{titlesec}
\usepackage{enumitem}
\usepackage{listings}
\usepackage[T1]{fontenc}
\usepackage{lmodern}
\usepackage{array}

\definecolor{codebg}{RGB}{248,248,248}
\definecolor{codeframe}{RGB}{210,210,210}
\definecolor{darkblue}{RGB}{20,55,100}

\lstdefinestyle{cppstyle}{
  language=C++,
  basicstyle=\ttfamily\footnotesize,
  backgroundcolor=\color{codebg},
  frame=single,
  rulecolor=\color{codeframe},
  columns=fullflexible,
  keepspaces=true,
  showstringspaces=false,
  breaklines=true,
  tabsize=4,
  xleftmargin=0.5em,
  xrightmargin=0.5em,
  framexleftmargin=0.4em,
  framexrightmargin=0.4em
}

\lstdefinestyle{bashstyle}{
  language=bash,
  basicstyle=\ttfamily\footnotesize,
  backgroundcolor=\color{codebg},
  frame=single,
  rulecolor=\color{codeframe},
  columns=fullflexible,
  keepspaces=true,
  showstringspaces=false,
  breaklines=true,
  tabsize=4,
  xleftmargin=0.5em,
  xrightmargin=0.5em,
  framexleftmargin=0.4em,
  framexrightmargin=0.4em
}

\titleformat{\section}{\large\bfseries\color{darkblue}}{}{0em}{}
\titleformat{\subsection}{\normalsize\bfseries}{}{0em}{}
\setlength{\parindent}{0pt}
\setlength{\parskip}{3pt}
\setlist[itemize]{leftmargin=1.2em,itemsep=1pt,topsep=2pt}

\begin{document}

\begin{center}
    {\LARGE \textbf{C++ Basics Cheat Sheet}}\\[2pt]
    {\small Direct syntax templates for C++ practice}
\end{center}

\hrule
\vspace{0.4em}

\begin{multicols}{2}

\section{1. Program skeleton}
\begin{lstlisting}[style=cppstyle]
#include <iostream>

int main() {
    // write program here
    return 0;
}
\end{lstlisting}

\section{2. Compile and run}
\begin{lstlisting}[style=bashstyle]
g++ -Wall -Wextra -std=c++17 path/to/file.cpp -o build/executable_name
./build/executable_name
\end{lstlisting}

\section{3. Basic types}
\begin{center}
\begin{tabular}{>{\ttfamily}l l}
\hline
int & whole number \\
double & decimal number \\
bool & true / false \\
char & one character \\
std::string & text \\
\hline
\end{tabular}
\end{center}

\begin{lstlisting}[style=cppstyle]
int count = 0;
double value = 2.5;
bool passed = true;
char label = 'A';
std::string name = "text";
\end{lstlisting}

\section{4. Constants}
\begin{lstlisting}[style=cppstyle]
const double fixed_value = 240.0;
double result = fixed_value*fixed_value;
\end{lstlisting}

\section{5. Printing}
\begin{lstlisting}[style=cppstyle]
std::cout << "label = " << variable << "\n";
std::cout << "a = " << a << ", b = " << b << "\n";
\end{lstlisting}

\section{6. If / else}
\begin{lstlisting}[style=cppstyle]
if (condition) {
    // code if true
} else if (other_condition) {
    // code if second condition is true
} else {
    // code otherwise
}
\end{lstlisting}

\section{7. Operators}
\begin{center}
\begin{tabular}{>{\ttfamily}c l}
\hline
= & assignment \\
== & equality check \\
!= & not equal \\
<, > & less / greater \\
<=, >= & less/greater or equal \\
\&\& & and \\
|| & or \\
! & not \\
\hline
\end{tabular}
\end{center}

\section{8. For loop}
Index loop:
\begin{lstlisting}[style=cppstyle]
for (int index = 0; index < number_of_steps; index++) {
    // repeated code
}
\end{lstlisting}

Range loop over a vector:
\begin{lstlisting}[style=cppstyle]
for (type element : vector_name) {
    // use element
}
\end{lstlisting}

\section{9. While loop}
Condition loop:
\begin{lstlisting}[style=cppstyle]
while (condition) {
    // repeated code
}
\end{lstlisting}

Read values from an input stream:
\begin{lstlisting}[style=cppstyle]
while (input_stream >> variable) {
    // variable stores one value read from the stream
}
\end{lstlisting}

Here \texttt{input\_stream} is the file/input object. \texttt{variable} stores the value read.

\section{10. Functions}
Template:
\begin{lstlisting}[style=cppstyle]
return_type function_name(input_type input_name) {
    // code
    return result;
}
\end{lstlisting}

Example form:
\begin{lstlisting}[style=cppstyle]
double function_name(double input_value) {
    return input_value*input_value;
}
\end{lstlisting}

\section{11. Structs}
Define a new grouped type:
\begin{lstlisting}[style=cppstyle]
struct TypeName {
    double member1;
    double member2;
};
\end{lstlisting}

Create and access:
\begin{lstlisting}[style=cppstyle]
TypeName object_name{value1, value2};
std::cout << object_name.member1 << "\n";
\end{lstlisting}

\section{12. Vectors}
\begin{lstlisting}[style=cppstyle]
#include <vector>

std::vector<double> values{1.0, 2.0, 3.0};
values.push_back(4.0);

for (double value : values) {
    std::cout << value << "\n";
}
\end{lstlisting}

Indexing:
\begin{lstlisting}[style=cppstyle]
values[0]  // first element
values[1]  // second element
\end{lstlisting}

\section{13. Strings}
\begin{lstlisting}[style=cppstyle]
#include <string>

std::string text_value = "some text";
std::cout << text_value << "\n";
\end{lstlisting}

\section{14. Read file}
\begin{lstlisting}[style=cppstyle]
#include <fstream>

std::ifstream input_file("path/to/input.txt");

if (!input_file) {
    std::cout << "Could not open input file\n";
    return 1;
}
\end{lstlisting}

Here \texttt{"path/to/input.txt"} is the file path. \texttt{input\_file} is the object used to read it.

\section{15. Write file}
\begin{lstlisting}[style=cppstyle]
#include <fstream>

std::ofstream output_file("path/to/output.txt");

if (!output_file) {
    std::cout << "Could not open output file\n";
    return 1;
}

output_file << "label = " << value << "\n";
\end{lstlisting}

Here \texttt{"path/to/output.txt"} is the output path. \texttt{output\_file} writes to it.

\section{16. Read lines}
\begin{lstlisting}[style=cppstyle]
#include <string>

std::string line;

while (std::getline(input_file, line)) {
    // line stores one full line from the file
}
\end{lstlisting}

\section{17. Parse text line}
\begin{lstlisting}[style=cppstyle]
#include <sstream>

std::stringstream line_stream(line);
std::string word;
double number = 0.0;

line_stream >> word >> number;
\end{lstlisting}

Here \texttt{line} is the full text line. \texttt{word} receives text. \texttt{number} receives a numeric value.

\section{18. Header file}
Header: \texttt{file\_name.hpp}
\begin{lstlisting}[style=cppstyle]
#pragma once

struct TypeName {
    double member;
};

return_type function_name(TypeName input);
\end{lstlisting}

Source: \texttt{file\_name.cpp}
\begin{lstlisting}[style=cppstyle]
#include "file_name.hpp"

return_type function_name(TypeName input) {
    // function definition
    return result;
}
\end{lstlisting}

\section{19. CMake template}
\begin{lstlisting}[style=cppstyle]
cmake_minimum_required(VERSION 3.16)
project(project_name)

set(CMAKE_CXX_STANDARD 17)
set(CMAKE_CXX_STANDARD_REQUIRED ON)

add_executable(executable_name
    main.cpp
    other_source.cpp
)
\end{lstlisting}

Build:
\begin{lstlisting}[style=bashstyle]
cmake -S . -B build              # configure project
cmake --build build              # compile project
./build/executable_name          # run executable
\end{lstlisting}

\section{20. Safe division}
\begin{lstlisting}[style=cppstyle]
double ratio = 0.0;

if (denominator > 0) {
    ratio = static_cast<double>(numerator) / denominator;
}
\end{lstlisting}

\section{21. Debugging order}
\begin{enumerate}
    \item Read the first compiler error.
    \item Check file name and line number.
    \item Fix one thing.
    \item Compile again.
\end{enumerate}

\section{22. Quick file checks}
\begin{lstlisting}[style=bashstyle]
pwd                              # show current folder
ls                               # list files here
cat file.cpp                     # print file content
grep -n "int main" file.cpp      # find main with line number
git status --short               # show changed files
\end{lstlisting}

\end{multicols}

\end{document}
